\subsection{Related Work}

We compare \aebench with similar recent works that use agents 
  for reproducibility, build agents specifically for artifact evaluation,
  or that build benchmarks to test both such agents.

\paragraph{\aebench vs Artisan}

Artisan~\cite{baek-ase26} is an agent for reproducing numerical results from
  software engineering papers by generating a standalone reproduction script.
It evaluates 60 tasks from 23 papers, and uses an execution-based output judge
  and an LLM method judge to verify that the script produces the expected
  results and does not simply copy precomputed values.
The method judge distinguishes copied results, last-mile reproduction and
  full reproduction, and accepts the latter two as successful.
This is quite related to \aebench, since it also considers how the results
  were produced, and not only whether the numbers are correct.

However, Artisan formulates reproduction as a code generation task, where
  the agent needs to produce a script for reproducing a specific table.
\aebench instead evaluates external coding agents operating an existing
  artifact, including environment setup, build, benchmark preparation and
  execution of the experiments.
An important difference is that Artisan accepts last-mile reproduction,
  which can skip some expensive experimental steps and use existing raw data.
In \aebench we want to check whether the agent performed the actual
  experiment required by the artifact, and not only obtained similar results.
We also want to study the changes agents make during the reproduction,
  including cases where fixing an execution error modifies the experimental
  procedure.
Artisan verifies the submitted script and its reproduction method, whereas
  our goal is also to examine whether the agent followed the documented
  procedure, and which deviations resulted in an invalid reproduction.

\paragraph{\aebench vs ArtifactCopilot}

ArtifactCopilot~\cite{wu-arxiv26} is an agent-based framework for automating
  artifact evaluation, including environment preparation, experiment execution
  and badge decisions.
Its latest version uses ArtifactGuide, a rubric that defines evaluation
  requirements, evidence levels and different types of reviewer interventions.
In particular, it distinguishes environment repairs from changes that modify
  the original experiment, which is also relevant to our work.
The system records execution evidence and checks whether the evidence supports
  the reviewer's conclusions, and it compares the resulting badge decisions
  against human evaluation on 60 artifacts.
\aebench has a different objective, since we evaluate the coding agents that
  perform artifact reproduction, rather than developing an agent to review
  and assign badges.
We use independent artifact-specific oracles to measure reproduction progress
  and record the execution of external coding agents.
This allows us to study how different agents react to failures, which changes
  they introduce in the artifact, and whether they continue with the original
  experiments or execute a modified procedure.
ArtifactCopilot already considers the validity of reviewer interventions,
  but does not systematically evaluate these behaviors of general coding agents
  under controlled reproduction tasks.
In particular, we are interested in cases where an agent passes the output
  checks even if the reproduced experiment was changed.

\paragraph{\aebench vs MLSys AE Agent}
Yang~\cite{yang-blog26} describes an artifact-evaluation agent deployed as a reviewer
  at MLSys\textquotesingle26, connected with HotCRP and a real review machine.
It reads papers, runs software, asks authors questions and prepares a badge decision,
  so it is not only a simple benchmark run.
It can return after an author response and continue the review, which makes its
  execution different from one isolated reproduction trial with a fixed input.
\aebench is not a conference reviewer and does not manage the author discussion.
This is a substantial system for actual peer review, including comments and badge
  decisions, whereas AEBench does not automatically become a reviewer.
Instead it provides repeatable cases and oracles to study what different agents do
  when reproducing the same artifact.

\paragraph{\aebench vs Heye et al.}
Heye et al.~\cite{heye-bigdata25} present an LLM-based toolkit for supporting
  artifact evaluation of security research papers.
Their approach consists of three stages: estimating the reproducibility of
  an artifact from its paper and documentation, preparing an execution
  environment using an LLM agent, and detecting methodological problems
  in the original research.
The agent executes shell commands inside a Docker container, installs
  dependencies and tries to resolve errors during execution.
They evaluate their approach against human reproducibility assessments,
  and their agent can prepare around 28\% of the artifacts previously
  found runnable by human reviewers.
This is close to \aebench, particularly the automated execution of
  artifact instructions in isolated environments.

However, Heye et al. mainly aim to reduce the manual work of reviewers,
  and measure whether their agent can prepare a runnable environment.
\aebench instead evaluates different coding agents on the reproduction
  process, using artifact-specific oracles for environment setup, build,
  benchmark preparation, and experiment execution.
Interestingly, Heye et al. also report an agent attempting to comment out
  an entire program to make it execute successfully, which would produce
  an incorrect reproducibility assessment.
They address this problem by modifying the agent prompts.
We are interested in studying this kind of behavior more systematically,
  including when agents change scripts, workloads, or configurations
  to overcome failures, and whether these changes invalidate the
  experiment even when the expected results are produced.

\paragraph{\aebench vs CORE-Bench}
CORE-Bench~\cite{siegel-tmlr25} evaluates agents that reproduce numerical results from
  provided research code and data, using 270 tasks from 90 papers.
The tasks have different difficulty levels and the agent must also find answers from
  the execution outputs.
\aebench instead uses artifact-specific checks for the environment, build, benchmark
  preparation, and experiments.
The distinction is that answering correctly about some results does not tell us
  exactly which steps the agent performed.

\paragraph{\aebench vs PaperBench}
PaperBench~\cite{starace-icml25} evaluates whether an agent can replicate the work of
  an AI paper, starting from its description and without using the original
  implementation.
It has 20 papers and a large hierarchy of grading rubrics, with an LLM judge to
  evaluate the work produced by agents.
This is a different task from \aebench, where the submitted artifact and its README
  are already given.
We are interested in operating that software correctly, and not in implementing a
  replacement.

\paragraph{\aebench vs REPRO-Bench}

REPRO-Bench~\cite{hu-acl25} evaluates whether AI agents can assess the
  reproducibility of social science papers using the original paper and
  its reproduction package.
It contains 112 papers with existing human reproduction reports, which
  provide a ground-truth reproducibility score from 1 to 4.
The agents inspect the code, execute reproduction scripts and compare
  the generated results with the findings reported in the paper.
The authors also analyze agent trajectories and identify several execution
  failures, including incorrect result comparisons, missing dependencies,
  and problems locating files in the reproduction package.
This is close to \aebench, as it evaluates different agents on real
  research artifacts, and also investigates why agents fail during
  reproduction.

However, REPRO-Bench measures whether the agent correctly predicts the
  reproducibility of a paper, using the human report as ground truth.
An agent may obtain the correct score even if it did not execute the
  reproduction procedure correctly.
\aebench instead evaluates the execution itself, using artifact-specific
  oracles for environment setup, build, benchmark preparation and
  experiment execution.
Also, REPRO-Bench analyzes agent trajectories to understand failures,
  but does not separately score whether the agent followed the
  documented procedure or introduced changes that invalidate it.
We are interested in these changes, particularly when agents modify
  scripts or experimental settings to overcome execution failures,
  and obtain the expected results using a different procedure.

\paragraph{\aebench vs DeployBench}
DeployBench~\cite{wang-arxiv26} is a benchmark of 51 tasks for deploying research
  artifacts from different domains onto an execution machine.
It uses hidden checks that execute a designated experiment and verify outputs, and
  studies when agents stop before the real task is completed.
\aebench shares the interest in environment setup and execution, but breaks the
  reproduction in four phases.
We also retain the execution evidence, since a working environment by itself is not
  the full reproduction.

\paragraph{\aebench vs Reward Hacking Benchmark}
The Reward Hacking Benchmark~\cite{thaman-icml26} gives tool-using agents tasks with
  opportunities to skip verification, use shortcuts or tamper with what gets checked.
It measures whether these shortcuts produce an undeserved success under different
  model configurations.
This is relevant to \aebench because agents may also change scripts or outputs to
  finish the reproduction.
However, our tasks come from real research artifacts, where a change can be a useful
  repair or an invalid experiment.

\paragraph{\aebench vs HackDetect}
HackDetect~\cite{shao-arxiv26} audits agent benchmark traces to identify shortcuts and
  explain when the obtained score does not represent the intended capability.
It studies protocol validity across multiple benchmarks and estimates the effect of
  exposed information on scores.
\aebench has a related concern, but the reference procedure is the artifact README and
  the expected research experiment.
Our checks also need to inspect the machine state, not only decide if a benchmark
  trace contains an exploit.